% !TeX root = ../papers/nc-p3.tex
\section{六行切断の余剰曲線と全被約窓での階数二曲面の排除}
\label{sec:rank-two-surface-excess-curves}
\begin{theorem}\label{thm:rank-two-surface-excess-curves}
被約な四辺有限の標準有限Frobenius窓で全階数反対平方根が一般可逆なら、階数二以下の軌跡は高々一次元である。
\end{theorem}
\begin{proof}[標準表示と四次元部分空間]
定理\ref{thm:rank-one-total-five}から階数二曲面があれば唯一の平面 $H=(h=0)$ である。
系\ref{cor:rank-two-plane-independent-curves}により $H$ は階数一曲線を含まない。
よって $H$ 上の一般階数二の二因子分解は定理\ref{thm:low-rank-surface-factors}で使った通り
二つの線型埋込みで与えられる。実際この分解に使った単連結性は $H$ からその上の有限個の階数一点を
除くだけでよい。二つの像平面は異なり、$H$ 上の比例点は高々二つである。特に平面外印数は $n\ge18$。
二つの一次列を $u,v$ とし、同定理の
\[
 A=\begin{pmatrix} C&u&v\\ u^t&0&-h\\ v^t&-h&0\end{pmatrix},
 \qquad G=\det A=-\det N/h^2
\]
を使う。全点で階数は二以上。平面外では階数が $2+\operatorname{rank}N$、平面上でも二列が独立なら
四以上である。$A$ の階数四以下の軌跡は高々一次元である。平面外で他の階数二曲面はなく、
$H$ 全体で四以下なら $h^2\mid G$ で、残余四次式が十八印以上を三重に通ることになる。
これは定理\ref{thm:low-rank-surface-factors}の四次式の重複度上界に反する。

一般の $W\subset k^6$、$\dim W=4$ を取り、
$\mathbb P^3\times\mathbb P(W)$ での六つの双一次式 $A(x)z=0$ の基点を考える。
階数五の曲面では核直線が $W$ に入る条件が余次元二、階数四の曲線では核二空間が $W$ と交わる条件が
余次元一なので、この二層からの部分は高々零次元である。これは $x$ も動かす入射の次元を数えることによる。
さらに階数四で二空間全体が $W$ に入る条件は余次元四で一般には起きない。
階数三の層は有限個の点と、階数一曲線の各成分の開部分だけで、各点の核三空間と $W$ は直線で交わる。
二次元以上になる条件の余次元は二だからである。階数二の点は $H$ 上の有限個で、
その核四空間との交わりは二次元である。従って基点の次元は高々一である。
\end{proof}
\begin{proof}[核直線が曲線上で使う次数]
階数一曲線の総被約次数を $d$ とすると $d\le2$。
$d=0$ は定理\ref{thm:low-rank-surface-exclusion}で処理済みなので $d>0$ とする。
各成分 $J$ は直線または滑らかな二次曲線であり $H$ に含まれない。$e=\deg J$ とする。
$J$ 上の出力平方方向の超平面次数も $e$ である。
$J$ の一般点での $A$ の核をGrassmann多様体の固有性で全曲線に延長し、
階数三の部分束 $K\subset k^6\otimes\mathcal O_J$ とする。
最初の四成分への射影は出力方向の零化超平面束 $T$ に入り、一般点で同型である。
これは $h\ne0$ の最後の二行の消去による。$T$ は自明四空間から出力の超平面束への
全射の核なので $\deg T=-e$。よって $-\deg K\ge e$ である。
一般 $W$ では曲線全体で $K\to(k^6/W)\otimes\mathcal O_J$ が全射で、核直線束の次数は $\deg K$。
従って基点に属するそのグラフの $(1,1)$ 次数は $e-\deg K\ge2e$ となる。

少なくとも一成分と $H$ との交点 $p$ を一つ取る。そこで $A$ の階数は三以下であり、
$u,v$ は比例する。階数三なら $K(p)=\ker A(p)$ に最後の二成分だけに支持された
非零ベクトルが入るから、$K\to T$ の行列式はそこで零になり、当該成分のグラフの次数が少なくとも一増す。
階数二なら $\mathbb P(W\cap\ker A(p))$ はグラフとは別の垂直直線である。
よって六式の基点に含まれる曲線の $(1,1)$ 総次数は少なくとも $2d+1$ である。
\end{proof}
\begin{proof}[連結な完全交叉からの残余]
五つの一般線型結合は純一次元の完全交叉 $Y$ を作る。実際基点の外で各段で一般切断が次元を下げ、
基点は既に一次元以下である。その $(1,1)$ 次数は $\binom63=20$ で、台は連結である。
例えば五切断のKoszul分解では、負の $( -j,-j)$ のコホモロジーは零か次数六のみにあり、
$j\le5$ より $H^0(\mathcal O_Y)=k$ となる。
第六の式で全体が消える成分と残りの成分を分ける。前者の被約次数だけでも $2d+1$ 以上である。
後者は空でない。階数一曲線外の平面外印が $n-3d>0$ 個以上あり、一般 $W$ でそれらは孤立基点になるからである。
連結性により両者の台は少なくとも一点で交わる。その点も第六の式で消えるが、孤立基点ではない。
従って孤立基点の相異なる個数は $20-(2d+1)-1$ 以下。ここでは残りの各被約成分に
第六の式を制限して次数で数えるだけなので、非被約な完全交叉でも同じ上界である。

$G$ の一般階数五の成分で核像が曲面 $Z\subset\mathbb P^5$ なら、その一般余次元二切断は
$\deg Z$ 個の新たな孤立基点を与える。前述の $n-3d$ との和を比較すると
\[
                 \deg Z\le 20-n+3d-(2d+2)\le20-n\le2 .
\]
一方、定理\ref{thm:low-rank-surface-exclusion}の最後の重複度評価は
十八以上の平面外印のみを使い、$G$ に次数五以上の既約成分が必要であることを示す。
この次数六の $G$ は従って被約である。その高次成分で核像のベクトル包の次元を $r$ とすると、
同定理の小行列式の議論で $r\ge5$。曲面像の次数は $r-2\ge3$ で上と矛盾する。
曲線像の場合の母線正規化と双対性による排除も同定理の証明を適用できる。
そこでの有限性には全点で $A$ の階数が二以上であることだけを使い、一般点では被約で階数五である。
これで最後の可能性も除かれる。
\end{proof}
